\section{Through the Weeks}

\subsection{The First Week}\label{sec:week1}

In Year B the first reading of the First Sunday is a prayer from the last
chapters of Isaiah (63:16b--17, 19b; 64:2--7), in which the prophet asks God
to tear the sky open:

\begin{quote}
O that thou wouldst rend the heavens, and wouldst come down: the mountains would
melt away at thy presence.\par
{\smallskip\noindent\small ---Isaiah 64:1 (63:19b in the Lectionary)\par}
\end{quote}

\noindent The Latin is \latin{Utinam dirumperes cælos, et descenderes}. The
prayer ends in trust: ``And now, O Lord, thou art our father, and we are clay:
and thou art our maker, and we all are the works of thy hands.'' The psalm of
the day is the eightieth (the seventy-ninth in the Douay), which prays:
``Convert us, O God: and shew us thy face, and we shall be saved.'' St Paul then thanks God that the
Corinthians lack no grace, ``waiting for the manifestation of our Lord Jesus
Christ,'' and the Gospel commands the watch (Section~\ref{sec:watch}). In Year A
Isaiah sees all nations flowing to the mountain of the Lord, where ``they shall
turn their swords into ploughshares,'' and St Paul sounds the hour to rise from
sleep, in the passage whose closing verses St Augustine read in the garden at
Milan on the day of his conversion. In Year C Jeremiah promises ``the bud of justice to
spring forth unto David,'' and St Luke
gives the signs in sun and moon and the command to lift up your heads.

The weekdays begin Isaiah in earnest. On Tuesday the prophet sees the rod out of
the root of Jesse; on Wednesday the feast of fat things on the mountain, when
the Lord will ``wipe away tears from every face,'' while the Gospel shows Christ
feeding the multitude on a mountain in Galilee; on Thursday the strong city,
with the Gospel of the house built on rock; on Friday the deaf who hear and the
blind who see, with the Gospel of the two blind men who follow Jesus crying
``Have mercy on us, O Son of David''; on Saturday the Teacher no longer hidden,
with the Gospel of the harvest and the labourers.

St Jerome, commenting on Tuesday's prophecy, read it of the Mother and the
Son: \latin{Nos autem virgam de radice Jesse sanctam Mariam Virginem
intelligamus, quæ nullum habuit sibi fruticem cohærentem}. Let us understand the
rod from the root of Jesse, he says, to be the holy Virgin Mary, who had no
other shoot growing with her; \latin{Et florem Dominum Salvatorem}, and the flower
to be the Lord our Saviour, who says in the Song, I am the flower of the field
and the lily of the valleys. On the next verse, where the Spirit of the Lord shall
rest upon him, Jerome quotes a gospel outside the canon, the Gospel
according to the Hebrews. In it the Spirit, descending on the Lord as he came up
from the water, says to him: \latin{Fili mi, in omnibus prophetis exspectabam
te, ut venires, et requiescerem in te}; my son, in all the prophets I was
waiting for thee, that thou mightest come and I might rest in thee. The words
are not Scripture, and the Church does not read them.

In the Office of Readings the
week's second readings come from St Cyril on Sunday, St Charles Borromeo's
pastoral letter on the season of Advent on Monday, St Gregory Nazianzen on
Tuesday, St Bernard on Wednesday, St Ephrem's commentary on the Diatessaron, an
early harmony of the Gospels, on Thursday, St Anselm on Friday and St Cyprian on Saturday.

The Gospel of Monday is the centurion of Capharnaum, who asked Jesus to heal his
servant and would not have him come to the house:

\begin{quote}
And the centurion, making answer, said: Lord, I am not worthy that thou shouldst
enter under my roof; but only say the word, and my servant shall be healed.\par
{\smallskip\noindent\small ---Matthew 8:8\par}
\end{quote}

\noindent St Augustine saw the paradox: ``By calling himself unworthy, he showed
himself worthy for Christ to come not into his house, but into his heart.'' An
old preacher whom the Catena Aurea cites under the name of Origen drew the
lesson that every Catholic knows:

\begin{quote}
And when you eat and drink the Lord's Body then the Lord enters under your roof,
and you then should humble yourself, saying, ``Lord, I am not worthy.''\par
{\smallskip\noindent\small ---Homily ascribed to Origen, in the Catena Aurea on Matthew 8:8\par}
\end{quote}

\noindent The Church put the centurion's words into the Mass with one change.
The Missal before the reform gave them so:
\latin{Domine, non sum dignus ut intres sub tectum meum; sed tantum dic verbo, et
sanabitur anima mea}, ``Lord, I am not worthy that thou shouldst enter under my
roof; say but the word, and my soul shall be healed.'' Where the centurion said
\latin{puer meus}, my servant, the communicant says \latin{anima mea}, my soul.
In the present Missal priest and people say the words together once, before
Communion.

When the First Sunday falls on 29 November, that Monday is 30 November, the
feast of St Andrew. The feast has its own readings, and its Gospel is the call
of Peter and Andrew by the sea of Galilee (Matthew 4:18--22). In its Office of
Readings the Church reads St John Chrysostom on another page, St John's account
of how Andrew, one of the two disciples of John the Baptist who first followed
the Lord, went at once to fetch his brother:

\begin{quote}
He findeth first his brother Simon and saith to him: We have found the Messias,
which is, being interpreted, the Christ.\par
{\smallskip\noindent\small ---John 1:41\par}
\end{quote}

\noindent Chrysostom, expounding these verses (John 1:41--42), heard Advent in
Andrew's words:

\begin{quote}
For this word, ``we have found,'' is the expression of a soul which travails for
His presence, and looks for His coming from above, and is made overjoyed when
the looked-for thing has happened, and hastens to impart to others the good
tidings.\par
{\smallskip\noindent\small ---St John Chrysostom, Homily 19 on St John\par}
\end{quote}

On 3 December the Church keeps the memorial of St Francis Xavier, who carried
the Gospel to India and Japan. From Cochin, on the last day of 1543, he wrote to
his brethren in Rome:

\begin{quote}
It often comes into my mind to go round all the Universities of Europe, and
especially that of Paris, crying out every where like a madman, and saying to
all the learned men there whose learning is so much greater than their charity,
`Ah! what a multitude of souls is through your fault shut out of heaven and
falling into hell!'\par
{\smallskip\noindent\small ---St Francis Xavier, letter from Cochin, 1543\par}
\end{quote}

\noindent He hoped that such men, once moved, would say from the bottom of their
hearts: ``Lord, here am I; send me whithersoever it shall please Thee, even to
India!''

On 4 December falls St John Damascene, who wrote of the title the Church gives
our Lady:

\begin{quote}
Hence it is with justice and truth that we call the holy Mary the Mother of God.
For this name embraces the whole mystery of the dispensation. For if she who
bore Him is the Mother of God, assuredly He Who was born of her is God and
likewise also man.\par
{\smallskip\noindent\small ---St John Damascene, On the Orthodox Faith, III.12\par}
\end{quote}

On 6 December the Church may keep St Nicholas, bishop of Myra. The Roman
Martyrology of 1902 told one wonder of him:
\latin{de quo inter plura miraculorum insignia illud memorabile fertur, quod
Imperatorem Constantinum ab interitu quorumdam se invocantium, longe constitutus,
ad misericordiam per visum monitis deflexit et minis}, ``of whom, among many
signal miracles, this memorable one is told: that, though far away, by warnings
and threats in a vision he turned the Emperor Constantine from the death of
certain men who called upon him, to mercy.'' When his day falls on a Sunday of
Advent, it is not kept.

On 7 December comes St Ambrose, bishop of Milan and Doctor of the Church. He is
kept on that day because it was the day of his ordination as bishop, and the
Martyrology entered him under that title: \latin{Mediolani ordinatio sancti Ambrosii Episcopi, et Ecclesiae
Doctoris}, ``At Milan, the ordination of Saint Ambrose, Bishop and Doctor of the
Church.'' He wrote a hymn for the Lord's birth that Pope St Celestine quoted in
430:

\begin{verse}\itshape
Veni Redemptor gentium,\\
Ostende partum Virginis,\\
Miretur omne sæculum,\\
Talis decet partus Deum.
\end{verse}

\noindent The \emph{Hymnal Noted} put it into English:

\begin{verse}
Come, Thou Redeemer of the earth,\\
Come, testify Thy Virgin-birth:\\
All lands admire, all times applaud:\\
Such is the birth that fits a God.
\end{verse}

\noindent Appendix~\ref{sec:prayers} gives the hymn in Latin and four of its
stanzas in English. Sing its first verse at home on Ambrose's day; by
Christmas the children will know it.
